\textbf{Problem 2 answer} \GradescopeComment{This should be on page 7 of your PDF.}

\begin{enumerate}[label=(\alph*)]

\item[(a)]
  Where is the first step that makes an assertion that is not necessarily true?

  \BLUE{
    In Lemma
    \REPLACEME                   % 2 or 3
    ,
    Step
    \REPLACEME                  % e.g. 5.2.2
    .
  }
  
\item[(b)]
  In one or two sentences, what precisely is wrong with the reasoning given in that step?

  \begin{blue}
    \REPLACEME
  \end{blue}

\item[(c)]
  Here is an input with $n=7$ on which the algorithm fails to return a local maximum:

  \BLUE{
    \[
      \begin{array}{ccc|c|ccc}
        0 & 0 & 0 &  0 & 0 & 0 &  0 \\
        0 & 0 &  0 & 0 & 0 &  0 & 0 \\
        0 & 0 & 0 & 0 & 0 & 0 & 0\\ \hline
        0 & 0 & 0 & 0 & 0 & 0 & 0\\ \hline
        0 & 0 & 0 & 0 & 0 & 0 & 0\\
        0 & 0 & 0 & 0 & 0 & 0 & 0 \\
        0 & 0 & 0 & 0 & 0 & 0 & 0
      \end{array}
    \]
  }

\item[]
  Here is what the algorithm does when given that input:
  \begin{blue}
    \begin{steps}
      \step \REPLACEME
      \step \REPLACEME
      
      \vdots
    \end{steps}
  \end{blue}
  
\end{enumerate}


\newpage


\paragraph{Problem \arabic{problem} collaboration and citations}~\GradescopeComment{This should be on page 8 (top) of your PDF.}

I collaborated on this problem with the following people:
\begin{blue}
  \begin{itemize}
  \item \REPLACEME                    % "none" if none
  \end{itemize}
\end{blue}

My answers use ideas from the following sources (other than the lecture notes and textbooks): 
\begin{blue}
  \begin{itemize}
  \item \REPLACEME                    % "none" if none
  \end{itemize}
\end{blue}



%%% Local Variables:
%%% mode: latex
%%% TeX-master: "homework_2"
%%% End:
